\documentclass[titlepage,a4paper]{jsarticle}
\usepackage{../../sty/import}% 各種パッケージインポート
\usepackage{../../sty/title}% タイトルページの変更

%% タイトルページの変数
% レポートタイトル
\title{情報技術分野の研究活動における生成AI利用の利点と倫理的問題}
% 提出日
\expdate{\today}
% 科目名
\subject{研究倫理}
% 分野
\class{情報経営システム工学分野}
% 学年
\grade{M1}
% 学籍番号
\mynumber{24336488}
% 記述者
\author{本間三暉}

\begin{document}
% titleページ作成
\maketitle
% 情報技術分野の研究活動において，生成AI（特に大規模言語モデル）を利用する場合の利点と倫理的問題を整理しなさい．
% 特に，論文執筆，プログラム作成，データ分析，査読，教育データ利用の観点から論じなさい．
\subsubsection*{情報技術分野の研究活動において，生成AI（特に大規模言語モデル）を利用する場合の利点と倫理的問題を整理しなさい．}
昨今，大規模言語モデルを始めとした生成AIの活用が注目されている．
情報技術分野の研究活動でも，うまく使えば作業の効率化や新しい発想の支援につながる可能性がある一方で，研究の信頼性や個人情報の扱いに関わる倫理的な問題も多く存在する．
そのため，生成AIを単なる便利な道具として使うだけでなく，どのような場面で有効であり，どのような危険があるのかを整理する必要がある．

生成AIは，情報技術分野の研究活動において，論文執筆，プログラム作成，データ分析，査読，教育データ利用などを効率化する有用な道具である．
論文執筆では，文章の構成整理や表現の改善，英文校正に役立ち，研究内容を分かりやすく伝える補助になる．
プログラム作成では，コード生成，エラー原因の調査，デバッグ，テストコード作成などに利用でき，研究に必要な実装作業を効率化できる．
また，データ分析では，分析方法の候補を示したり，分析コードやグラフ作成を支援したりできる．
査読においても，論文の要点整理や確認すべき観点の洗い出しに利用できる可能性がある．
さらに，教育データの利用においては，学生一人ひとりの理解度や学習状況に合わせたフィードバックや教材作成に活用できる可能性がある．

一方で，生成AIの利用には倫理的な問題もある．論文執筆では，存在しない文献を生成したり，先行研究の内容を誤って説明したりする可能性がある．
そのような内容を確認せずに使うと，研究の根拠が不正確になり，論文全体の信頼性を損なう．
また，AIが作った文章を自分の考えとして提出すると，どこまでが研究者本人の考察なのかが曖昧になり，著者としての責任の所在も分かりにくくなる．

プログラム作成では，生成されたコードにバグや脆弱性が含まれる可能性がある．
特に研究で使うプログラムの場合，コードの誤りは実験結果や分析結果に直接影響する．
そのため，動いたから正しいと考えるのではなく，処理内容を理解し，テストや検証を行ったうえで利用する必要がある．
また，生成されたコードが既存のコードと似ている場合には，著作権やライセンス上の問題が生じる可能性もある．

データ分析では，不適切な分析手法や誤った解釈を示す可能性がある．
たとえば，データの性質に合わない統計手法を提案したり，相関関係を因果関係のように説明したりすることがある．
このような出力をそのまま使うと，研究結果の解釈を誤る危険がある．また，個人情報や未公開データを外部の生成AIに入力すると，情報漏洩やプライバシー侵害につながる可能性がある．

査読では，未公開論文を外部AIに入力することで機密性が損なわれるおそれがある．
査読対象の論文には，まだ公開されていない研究成果やアイデアが含まれているため，その内容を不用意に外部サービスへ入力することは問題である．
また，AIに評価を任せてしまうと，査読者自身の専門的な判断や責任が弱くなる．査読は研究者の知識と責任に基づいて行われるべきであり，生成AIはあくまで補助的に扱う必要がある．

教育データ利用では，個人情報や学習履歴の扱いを誤ると，学生のプライバシーや公平性に関わる問題が生じる．
成績，課題の提出状況，学習履歴などは学生個人に関わる重要な情報であり，不適切に利用されると本人に不利益を与える可能性がある．
また，AIによる評価や推薦に偏りがある場合，特定の学生に不公平な扱いが生じる危険もある．
そのため，教育データを利用する場合には，利用目的を明確にし，必要に応じて同意を得ることや，個人が特定されないようにデータを扱うことが重要である．

したがって，生成AIは研究者の代わりではなく，あくまで補助的な道具として利用すべきである．
生成AIの出力をそのまま信用せず，人間が内容を確認し，必要に応じて利用を明示し，最終的な責任を研究者自身が負うことが重要である．
また，未公開の研究内容や個人情報を安易に入力しないこと，生成された文章，コード，分析結果を必ず検証することも必要である．
生成AIを適切に使えば研究活動の効率化に役立つが，その利用には透明性，責任，プライバシー保護への配慮が不可欠である．



% % 参考文献
% \begin{thebibliography}{99}
% \bibitem{} 
% \end{thebibliography}

\end{document}